\chapter{從 C 到 C++}
\label{ch:from-c-to-cpp}

學完 C，你已經會變數、迴圈、陣列、函式、指標和結構。 C++ 幾乎把這些
全部保留下來，再往上加了參考、類別、模板和一整套標準函式庫。 換句話說，
你不是從零開始學一個新語言，而是在熟悉的地基上多蓋幾層。

這一章只做最基本的轉換：寫出第一支 C++ 程式並編譯它，看懂到處出現的
\texttt{std::}，學會用 \texttt{cout} 輸出、用 \texttt{cin} 輸入，以及
解題時常見的輸入加速寫法。 最後一節整理 C 的寫法在 C++ 裡哪些照用、
哪些要改。

\section{第一支 C++ 程式}
\label{sec:first-program}

先看熟悉的 C 版本：

\begin{code}{c}{C 的 Hello}
#include <stdio.h>

int main(void) {
    printf("Hello, C!\n");
    return 0;
}
\end{code}

對應的 C++ 版本長這樣：

\begin{example}[label=ex:hello]{cpp}{(hello.cpp) 第一支 C++ 程式}
#include <iostream>

int main() {
    std::cout << "Hello, C++!" << '\n';
    std::cout << 1 + 2 << '\n';
    return 0;
}
\end{example}

\iosample{}{
Hello, C++! \\
3 \\
}

逐行對照：

\begin{itemize}
\item \texttt{\#include <iostream>}：C++ 的輸入輸出放在
\texttt{<iostream>}（input/output stream）這個標頭檔裡，地位相當於
C 的 \texttt{<stdio.h>}。 C++ 標準函式庫的標頭檔沒有 \texttt{.h}。
\item \texttt{int main()}：C++ 裡空括號就代表沒有參數，不必寫
\texttt{void}（寫了也合法）。
\item \texttt{std::cout}：標準輸出，也就是螢幕。 前面的 \texttt{std::}
下一節解釋。
\item \texttt{<<}：把右邊的東西送進 \texttt{cout}。 可以一直串下去，
一行送好幾個值。
\item \texttt{'\textbackslash n'}：換行字元，跟 C 一樣。
\end{itemize}

第二行輸出值得注意：\texttt{1 + 2} 印出 $3$，我們沒有寫任何
\texttt{\%d}。 \texttt{cout} 會依照值的型別決定怎麼印，這是它跟
\texttt{printf} 最大的不同，第~\ref{sec:cout}~節再細看。

\subsection{編譯與執行}

C++ 原始檔的副檔名慣例是 \texttt{.cpp}。 編譯器改用 \texttt{g++}，
用法跟 \texttt{gcc} 幾乎一樣：

\begin{code}{bash}{編譯並執行 hello.cpp}
g++ -std=c++17 -Wall hello.cpp -o hello
./hello
\end{code}

\texttt{-std=c++17} 指定用 C++17 標準編譯，\texttt{-Wall} 打開常用警告。
本書每一支完整範例都用這行指令編譯、實際跑過，印出的結果就是書上寫的。

Windows 上執行檔叫 \texttt{hello.exe}，在命令提示字元裡直接打
\texttt{hello} 執行。 用 IDE 的話，記得在設定裡把語言標準選成 C++17
或更新的版本。

\subsection{常見錯誤}

\begin{itemize}
\item \textbf{用 \texttt{gcc} 編譯 \texttt{.cpp}}：\texttt{gcc} 看到
\texttt{.cpp} 會用 C++ 模式編譯，但連結時不會自動加上 C++ 標準函式庫，
常見的結果是一長串 \texttt{undefined reference}。 編 C++ 一律用
\texttt{g++}。
\item \textbf{寫成 \texttt{<iostream.h>}}：這是很舊的寫法，現在的編譯器
找不到這個檔案。 C++ 標準標頭檔沒有 \texttt{.h}。
\item \textbf{忘了 \texttt{-std=c++17}}：本書後面會用到 C++17 才有的
寫法。 舊版編譯器的預設標準可能比較舊，明確寫出來最保險。
\end{itemize}

\section{namespace 與 \texttt{std::}}
\label{sec:namespace}

大型程式和函式庫裡名字很多，難免撞名：你寫了一個 \texttt{count}，某個
函式庫也有一個 \texttt{count}。 C 只能靠取名字時加前綴來避開，例如
\texttt{mylib\_count}。 C++ 提供\textbf{命名空間}（namespace），把名字
分裝在不同的空間裡，用 \texttt{空間名::名字} 指定要哪一個。
\texttt{::} 叫做\textbf{範圍解析運算子}（scope resolution operator）。

\begin{example}[label=ex:namespaces]{cpp}{(namespaces.cpp) 兩個空間裡的同名變數}
#include <iostream>

namespace school {
    int count = 30;
}

namespace club {
    int count = 12;
}

int main() {
    std::cout << school::count << ' ' << club::count << '\n';
    return 0;
}
\end{example}

\iosample{}{
30 12 \\
}

兩個 \texttt{count} 住在不同的空間，互不干擾。 C++ 標準函式庫的名字
幾乎都放在一個叫 \texttt{std}（standard）的命名空間裡，所以要寫
\texttt{std::cout}，意思是 \texttt{std} 裡的 \texttt{cout}。

\subsection{\texttt{using namespace std;}}

每次都寫 \texttt{std::} 有點囉嗦。 在檔案開頭加一行

\begin{code}{cpp}{把 std 的名字全部引進來}
using namespace std;
\end{code}

之後就能直接寫 \texttt{cout}，編譯器會自己到 \texttt{std} 裡找。 解題
和小程式常這樣寫，本書從下一節開始也這樣寫，讓範例短一點。

代價是 \texttt{std} 裡上千個名字全部倒進來，撞名的機會變大。 下面這段
程式自己定義了全域變數 \texttt{count}，而 \texttt{<algorithm>} 裡剛好
有一個函式 \texttt{std::count}：

\begin{code}{cpp}{using namespace std 造成撞名（錯誤示範）}
#include <algorithm>
#include <iostream>
using namespace std;

int count = 0;

int main() {
    count++;               // 編譯錯誤
    cout << count << '\n';
    return 0;
}
\end{code}

\texttt{g++} 會說：

\begin{code}{text}{編譯器的錯誤訊息}
error: reference to 'count' is ambiguous
\end{code}

ambiguous 是有歧義：編譯器不知道你要自己的 \texttt{count} 還是
\texttt{std::count}。 遇到這種錯誤，把自己的變數改個名字就好。 寫給別
人用的程式或多人合作的專案，通常不寫 \texttt{using namespace std;}，
老實寫 \texttt{std::}。 折衷的寫法是只引進用得到的名字，例如
\texttt{using std::cout;}。

\subsection{常見錯誤}

\begin{itemize}
\item \textbf{漏了 \texttt{std::}}：沒有寫 \texttt{using}，又直接寫
\texttt{cout}，\texttt{g++} 會報 \texttt{'cout' was not declared in
this scope; did you mean 'std::cout'?}。 照提示補上 \texttt{std::}，
或在開頭加 \texttt{using}。
\item \textbf{寫在標頭檔裡}：\texttt{using namespace std;} 若寫在
\texttt{.h} 標頭檔，每個 \texttt{\#include} 它的檔案都會被迫引進整個
\texttt{std}。 標頭檔裡一律寫完整的 \texttt{std::}。
\end{itemize}

\section{用 \texttt{cout} 輸出}
\label{sec:cout}

\begin{crecap}[printf 的格式字串]
\texttt{printf} 靠格式字串決定怎麼印每個引數：\texttt{\%d} 印
\texttt{int}、\texttt{\%lld} 印 \texttt{long long}、\texttt{\%f} 印
\texttt{double}、\texttt{\%c} 印字元、\texttt{\%s} 印字串，
\texttt{\%.2f} 印到小數點後兩位。 格式與型別對不上（例如用
\texttt{\%d} 印 \texttt{double}）是未定義行為，編譯器只會給警告，
程式照樣編得過、印出錯的東西。
\end{crecap}

\texttt{cout} 不用格式字串。 每個用 \texttt{<<} 送進去的值，編譯器都
知道它的型別，會自動選對應的印法，所以不會有格式與型別對不上的問題。

\begin{example}[label=ex:cout-basics]{cpp}{(cout\_basics.cpp) 不同型別的輸出與小數位數}
#include <iomanip>
#include <iostream>
using namespace std;

int main() {
    int n = 3;
    double avg = 247.0 / n;
    char grade = 'B';
    bool passed = avg >= 60;

    cout << "人數：" << n << '\n';
    cout << "平均：" << avg << '\n';
    cout << "等第：" << grade << '\n';
    cout << "及格：" << passed << '\n';

    cout << fixed << setprecision(2);
    cout << "平均：" << avg << '\n';
    cout << "再印一次：" << 1.0 / 3 << '\n';
    return 0;
}
\end{example}

\iosample{}{
人數：3 \\
平均：82.3333 \\
等第：B \\
及格：1 \\
平均：82.33 \\
再印一次：0.33 \\
}

幾個要注意的地方：

\begin{itemize}
\item \texttt{double} 預設印 $6$ 位\textbf{有效數字}，所以
$82.333\ldots$ 印成 \texttt{82.3333}，不是 C 的 \texttt{\%f} 那種固定
小數點後 $6$ 位。
\item \texttt{bool} 是 C++ 內建的型別，值是 \texttt{true} 或
\texttt{false}，用 \texttt{cout} 印出來是 \texttt{1} 或 \texttt{0}。
\item 要固定小數位數，先送 \texttt{fixed}，再送
\texttt{setprecision(位數)}，效果等於 \texttt{\%.2f}。
\texttt{setprecision} 在 \texttt{<iomanip>} 裡，要另外 include。
\item \texttt{fixed} 和 \texttt{setprecision} 一旦設定就一直有效。
範例最後一行沒有再設定，\texttt{1.0 / 3} 照樣印成兩位小數。
\end{itemize}

\subsection{\texttt{endl} 與 \texttt{'\textbackslash n'}}

很多教材用 \texttt{endl} 換行：

\begin{code}{cpp}{兩種換行寫法}
cout << "Hello" << endl;   // 換行，並且立刻清空輸出緩衝區
cout << "Hello" << '\n';   // 只換行
\end{code}

兩者都會換行，差別在 \texttt{endl} 每次還會\textbf{清空緩衝區}
（flush），把累積的輸出立刻寫出去。 偶爾印幾行看不出差別；在迴圈裡印
幾十萬行時，每行都 flush 會明顯拖慢程式。 本書一律用
\texttt{'\textbackslash n'}。

\subsection{常見錯誤}

\begin{itemize}
\item \textbf{字元做運算後變成整數}：\texttt{char c = 'A';} 之後
\texttt{cout << c + 1} 印出 \texttt{66}，不是 \texttt{B}。
\texttt{c + 1} 的型別已經提升成 \texttt{int}，\texttt{cout} 就照整數
印。 要印字元，寫 \texttt{(char)(c + 1)}。
\item \textbf{忘了加括號}：\texttt{<<} 的優先順序比比較運算子高，
\texttt{cout << a < b;} 會被解讀成 \texttt{(cout << a) < b}，編譯失敗。
要印比較結果，寫 \texttt{cout << (a < b);}。
\item \textbf{只寫 \texttt{setprecision} 沒寫 \texttt{fixed}}：這時
\texttt{setprecision(2)} 指的是\textbf{有效數字}兩位，$82.3333$ 會印成
\texttt{82}。 要固定小數點後幾位，兩個都要寫。
\end{itemize}

\section{用 \texttt{cin} 輸入}
\label{sec:cin}

\begin{crecap}[scanf 為什麼要 \&]
C 的函式參數是傳值：函式拿到的是引數的複本，改複本不會影響呼叫者的變
數。 \texttt{scanf} 要把讀到的值寫回你的變數，只能請你把變數的位址交
給它，它再沿著指標寫進去：

\begin{code}{c}{scanf 收的是位址}
int n;
scanf("%d", &n);   // 傳 n 的位址，scanf 透過指標寫入 n
\end{code}

漏了 \texttt{\&}，\texttt{scanf} 會把 \texttt{n} 裡的垃圾值當成位址去
寫，結果是未定義行為。 另外，\texttt{scanf} 的回傳值是成功讀到幾個
項目，讀到檔案結尾時回傳 \texttt{EOF}。
\end{crecap}

C++ 用 \texttt{cin}（標準輸入）和 \texttt{>>} 讀資料，箭頭方向是從
\texttt{cin} 流向變數：

\begin{code}{cpp}{cin 的基本用法}
int n;
cin >> n;          // 不用寫 &
int a, b;
cin >> a >> b;     // 一次讀兩個
\end{code}

\texttt{cin >> n} 不用寫 \texttt{\&n}，卻一樣能改到 \texttt{n}。 背後
的機制叫\textbf{參考}（reference），是第~2~章的主題；現在只要記得
\texttt{cin} 後面直接寫變數名稱。

\texttt{>>} 會先跳過空白、Tab 和換行，再讀一個值，所以數字之間用空白
還是換行隔開都可以，這點跟 \texttt{scanf("\%d")} 一樣。

\begin{example}[label=ex:cin-sum]{cpp}{(cin\_sum.cpp) 讀 n 個整數，輸出總和與最大值}
#include <iostream>
using namespace std;

int main() {
    int n;
    cin >> n;

    int sum = 0, mx = 0;
    for (int i = 0; i < n; i++) {
        int x;
        cin >> x;
        sum += x;
        if (i == 0 || x > mx) {
            mx = x;
        }
    }
    cout << sum << ' ' << mx << '\n';
    return 0;
}
\end{example}

\iosample{
5 \\
3 9 -2 7 4 \\
}{
21 9 \\
}

\subsection{讀到檔案結尾為止}

有些題目不先告訴你有幾筆資料，要一直讀到輸入結束。 C 的寫法是
\texttt{while (scanf("\%d", \&x) == 1)}。 C++ 的寫法更短：

\begin{example}[label=ex:cin-eof]{cpp}{(cin\_eof.cpp) 讀到輸入結束，輸出筆數與總和}
#include <iostream>
using namespace std;

int main() {
    int x;
    int cnt = 0;
    long long sum = 0;
    while (cin >> x) {
        cnt++;
        sum += x;
    }
    cout << cnt << ' ' << sum << '\n';
    return 0;
}
\end{example}

\iosample{
10 20 30 \\
40 \\
}{
4 100 \\
}

\texttt{cin >> x} 這個運算式本身可以當條件用：讀取成功就是真，讀到結尾
或讀取失敗就是假，迴圈自然結束。 在終端機手動輸入時，Linux 和 macOS
按 Ctrl+D、Windows 按 Ctrl+Z 再按 Enter，代表輸入結束。

\subsection{常見錯誤}

\begin{itemize}
\item \textbf{箭頭方向寫反}：輸出是 \texttt{cout <<}，輸入是
\texttt{cin >>}。 箭頭指向資料要去的地方。 寫成 \texttt{cin << x} 會
編譯失敗。
\item \textbf{照 C 的習慣寫 \texttt{cin >> \&x}}：\texttt{\&x} 是位址，
不能拿來存讀進來的值，編譯會失敗。 這類錯誤的訊息常常有幾十行，先看第一
個 \texttt{error:} 所在的行號，通常就是寫錯的那一行。
\item \textbf{輸入格式不符就繼續讀}：要讀整數卻碰到 \texttt{abc}，這次
讀取失敗，\texttt{cin} 進入失敗狀態。 呼叫 \texttt{cin.clear()} 清除
這個狀態之前，後面的 \texttt{>>} 都不會再讀任何東西。 用 \texttt{while (cin >> x)} 這種寫法，失敗時迴圈會停下來，
比較不會出事。
\end{itemize}

\section{輸入加速}
\label{sec:fast-io}

解題時資料量可能很大，例如一次讀進幾十萬個整數。 這時 \texttt{cin}
和 \texttt{cout} 的預設設定會拖慢速度，常見的做法是在 \texttt{main} 一
開頭加兩行：

\begin{code}{cpp}{輸入加速的兩行}
ios::sync_with_stdio(false);
cin.tie(nullptr);
\end{code}

這兩行各自關掉一項預設行為：

\begin{itemize}
\item \texttt{ios::sync\_with\_stdio(false)}：預設情況下，
\texttt{cin}／\texttt{cout} 和 C 的 \texttt{scanf}／\texttt{printf}
保持同步，混著用時順序不會亂。 同步有成本，這行把它關掉。
\item \texttt{cin.tie(nullptr)}：預設情況下 \texttt{cin} 綁在
\texttt{cout} 上，每次讀輸入前都會先 flush \texttt{cout}，確保提示文字
已經印出來。 這行解除綁定。 \texttt{nullptr} 是 C++ 的空指標，相當於
C 的 \texttt{NULL}，第~\ref{sec:c-in-cpp}~節會再提。
\end{itemize}

\begin{example}[label=ex:fastio]{cpp}{(fastio.cpp) 加速輸入輸出，逐行印出前綴和}
#include <iostream>
using namespace std;

int main() {
    ios::sync_with_stdio(false);
    cin.tie(nullptr);

    int n;
    cin >> n;
    long long prefix = 0;
    for (int i = 0; i < n; i++) {
        int x;
        cin >> x;
        prefix += x;
        cout << prefix << '\n';
    }
    return 0;
}
\end{example}

\iosample{
5 \\
3 9 -2 7 4 \\
}{
3 \\
12 \\
10 \\
17 \\
21 \\
}

輸出用的是 \texttt{'\textbackslash n'} 而不是 \texttt{endl}。 前面那兩
行省下來的時間，如果每一行都 \texttt{endl} 一次，很容易又被吃回去。

\subsection{常見錯誤}

關掉同步之後，\texttt{cout} 和 \texttt{printf} 各自有自己的緩衝區，
誰先寫到螢幕上就不再有保證。 下面這段程式在本書測試的環境下，
\texttt{printf} 那行跑到最後面：

\begin{code}{cpp}{關掉同步後混用 printf（錯誤示範）}
ios::sync_with_stdio(false);
cout << "第一行" << '\n';
printf("第二行\n");
cout << "第三行" << '\n';
\end{code}

\iosample{}{
第一行 \\
第三行 \\
第二行 \\
}

換一個環境，順序可能又不一樣。 規則很簡單：加了這兩行，整支程式就只
用 \texttt{cin}／\texttt{cout}，不要再混用 \texttt{scanf}／
\texttt{printf}。

另一個常見錯誤是把這兩行放在已經讀過輸入之後。 它們要在第一次輸入輸出
之前執行，所以放在 \texttt{main} 的最前面。

\section{C 的寫法在 C++ 裡}
\label{sec:c-in-cpp}

C++ 幾乎是 C 的超集合：你在 C 學過的語法，絕大部分在 C++ 裡照樣能
用。 這一節整理需要留意的幾個地方。

\subsection{C 的標頭檔}

C 標準函式庫在 C++ 裡都還在，只是標頭檔換個名字：去掉 \texttt{.h}，前
面加 \texttt{c}。

\begin{center}
\begin{tabular}{ll}
\toprule
C & C++ \\
\midrule
\texttt{<stdio.h>} & \texttt{<cstdio>} \\
\texttt{<stdlib.h>} & \texttt{<cstdlib>} \\
\texttt{<math.h>} & \texttt{<cmath>} \\
\texttt{<string.h>} & \texttt{<cstring>} \\
\texttt{<ctype.h>} & \texttt{<cctype>} \\
\bottomrule
\end{tabular}
\end{center}

\texttt{printf}、\texttt{sqrt}、\texttt{strlen} 這些函式的用法都跟 C
一樣。 寫成 \texttt{<stdio.h>} 通常也編得過，不過 C++ 程式慣用
\texttt{c} 開頭的版本。

\subsection{\texttt{bool} 是內建型別}

C 要 \texttt{\#include <stdbool.h>} 才有 \texttt{bool}、\texttt{true}、
\texttt{false}。 在 C++ 裡它們是語言本身的一部分，不用 include。

\subsection{空指標寫 \texttt{nullptr}}

\begin{crecap}[NULL 與 malloc]
指標不指向任何東西時設成 \texttt{NULL}，使用前先檢查。
\texttt{malloc} 向 heap 要一塊記憶體，回傳型別是 \texttt{void *}，用完
要 \texttt{free}：

\begin{code}{c}{C 的動態配置}
int *a = malloc(5 * sizeof(int));   // void * 自動轉成 int *
if (a != NULL) {
    /* 使用 a[0] .. a[4] */
    free(a);
}
\end{code}

C 允許 \texttt{void *} 自動轉成任何物件指標，所以不必轉型。
\end{crecap}

C++ 有兩處不同：

\begin{itemize}
\item 空指標建議寫 \texttt{nullptr}。 它是專門給指標用的空值，不會像
\texttt{NULL} 那樣在某些地方被當成整數 $0$。
\item C++ 不允許 \texttt{void *} 自動轉成 \texttt{int *}。 上面那行
\texttt{malloc} 原封不動放進 \texttt{.cpp}，\texttt{g++} 會報
\texttt{invalid conversion from 'void*' to 'int*'}，要自己加上
\texttt{(int *)} 轉型。
\end{itemize}

其實寫 C++ 時很少需要 \texttt{malloc}。 要一段長度可變的陣列，第~3~章
的 \texttt{vector} 會幫你管記憶體；真的要自己配置，C++ 有
\texttt{new}／\texttt{delete}，第~4~章再談。 下面的範例只是示範 C 的
寫法在 C++ 裡還能跑：

\begin{example}[label=ex:c-in-cpp]{cpp}{(c\_in\_cpp.cpp) 在 C++ 裡用 C 的函式庫}
#include <cmath>
#include <cstdio>
#include <cstdlib>
#include <iostream>
using namespace std;

bool isSquare(int n) {
    int r = (int)sqrt(n);
    return r * r == n;
}

int main() {
    int n = 5;
    int *a = (int *)malloc(n * sizeof(int));
    for (int i = 0; i < n; i++) {
        a[i] = i * i + 9;
    }
    for (int i = 0; i < n; i++) {
        printf("%d:%s ", a[i], isSquare(a[i]) ? "yes" : "no");
    }
    printf("\n");
    free(a);

    int *p = nullptr;
    if (p == nullptr) {
        cout << "p 沒有指向任何東西" << '\n';
    }
    return 0;
}
\end{example}

\iosample{}{
9:yes 10:no 13:no 18:no 25:yes \\
p 沒有指向任何東西 \\
}

這支程式沒有關掉同步，所以 \texttt{printf} 和 \texttt{cout} 混用，順序
仍然正確。

\subsection{常見錯誤}

\begin{itemize}
\item \textbf{把變數取名為 C++ 的關鍵字}：C++ 多了一些 C 沒有的關鍵字，
例如 \texttt{new}、\texttt{delete}、\texttt{class}、\texttt{this}、
\texttt{template}。 在 C 裡把變數取名為 \texttt{new} 沒問題，搬到 C++
就編不過，要改名。
\item \textbf{以為 C++ 會幫你檢查陣列越界}：C 的陣列在 C++ 裡還是原來
的陣列，越界一樣是未定義行為，編譯器和執行時都不一定會提醒你。 第~3~
章的 \texttt{vector} 提供有檢查的存取方式。
\end{itemize}

\section{本章練習}
\label{sec:ch1-exercises}

\begin{drills}[label=drills:ch1]
\item 預測下面程式的輸出：
\begin{code}{cpp}{練習：預測輸出}
#include <iomanip>
#include <iostream>
using namespace std;

int main() {
    cout << fixed << setprecision(1) << 3.14159 << '\n';
    cout << 2.0 << '\n';
    cout << 'A' + 1 << ' ' << (char)('A' + 1) << '\n';
    return 0;
}
\end{code}
\item 下面的程式編譯失敗，找出原因並用兩種不同的方法修好：
\begin{code}{cpp}{練習：找錯}
#include <iostream>

int main() {
    int n;
    cin >> n;
    cout << n * 2 << '\n';
    return 0;
}
\end{code}
\item 寫一支程式：讀入兩個整數 $a$、$b$（$b \neq 0$），分三行輸出
$a + b$、$a \times b$，以及 $a / b$ 的實數結果（小數點後三位）。 輸入
\texttt{7 2} 時應輸出 \texttt{9}、\texttt{14}、\texttt{3.500}。
\item 一支程式開頭加了 \texttt{ios::sync\_with\_stdio(false);}，後面
同時用了 \texttt{cout} 和 \texttt{printf}。 可能會出什麼問題？ 該怎麼
改？
\item 寫一支程式：不斷讀入整數直到輸入結束，輸出讀到幾個數，以及它們的
平均（小數點後兩位），兩者以空白隔開。 輸入 \texttt{3 4 8} 時應輸出
\texttt{3 5.00}。 讀不到任何數時輸出 \texttt{0}。
\item 為什麼大量輸出時建議用 \texttt{'\textbackslash n'} 而不是
\texttt{endl}？
\end{drills}

\addcontentsline{toc}{section}{練習解答}
\section*{練習解答}

\begin{enumerate}
\item 輸出三行：\texttt{3.1}、\texttt{2.0}、\texttt{66 B}。
\texttt{fixed} 和 \texttt{setprecision(1)} 設定後一直有效，所以
\texttt{2.0} 也印一位小數；\texttt{'A' + 1} 是 \texttt{int}，印出
$66$，轉回 \texttt{char} 才印 \texttt{B}。
\item \texttt{cin}、\texttt{cout} 在 \texttt{std} 裡，沒寫
\texttt{std::}，編譯器找不到。 改法一：寫成 \texttt{std::cin}、
\texttt{std::cout}。 改法二：在 \texttt{\#include} 下面加
\texttt{using namespace std;}（或 \texttt{using std::cin;} 與
\texttt{using std::cout;}）。
\item 一種寫法：
\begin{code}{cpp}{練習 3 參考解答}
#include <iomanip>
#include <iostream>
using namespace std;

int main() {
    int a, b;
    cin >> a >> b;
    cout << a + b << '\n';
    cout << a * b << '\n';
    cout << fixed << setprecision(3) << (double)a / b << '\n';
    return 0;
}
\end{code}
要先把 \texttt{a} 轉成 \texttt{double} 再除，否則 \texttt{7 / 2} 是整
數除法，得到 $3$。
\item 關掉同步後，\texttt{cout} 和 \texttt{printf} 各用各的緩衝區，輸出
順序可能跟程式碼順序不同。 改法是整支程式只用 \texttt{cout}（或者拿掉
那一行，保留同步）。
\item 一種寫法：
\begin{code}{cpp}{練習 5 參考解答}
#include <iomanip>
#include <iostream>
using namespace std;

int main() {
    ios::sync_with_stdio(false);
    cin.tie(nullptr);

    int x, cnt = 0;
    long long sum = 0;
    while (cin >> x) {
        cnt++;
        sum += x;
    }
    if (cnt == 0) {
        cout << 0 << '\n';
        return 0;
    }
    cout << cnt << ' ' << fixed << setprecision(2) << (double)sum / cnt << '\n';
    return 0;
}
\end{code}
\item \texttt{endl} 除了換行，每次還會 flush 緩衝區；印很多行時，每行都
flush 會拖慢程式。 \texttt{'\textbackslash n'} 只換行，讓輸出累積在緩衝
區裡一次寫出去。
\end{enumerate}
